\documentclass[../../main.tex]{subfiles}% 注意这里的文件路径不能用 ./main.tex ,否则用latexmk编译子文件会报错
\graphicspath{{\subfix{./image/}}} % 指定图片目录,后续可以直接使用图片文件名
% 注意这里的文件路径不能用 ../../image/ ,否则用latexmk编译子文件会报错

% 例如:
% \begin{figure}[H]
% \centering
% \includegraphics[scale=0.3]{图.png}
% \caption{}
% \label{figure:图}
% \end{figure}
% 注意:上述\label{}一定要放在\caption{}之后,否则引用图片序号会只会显示??.

\begin{document}

\section{基本定理}

常见的反例:$f(x)=x^m \sin\frac{1}{x^n}.$

\begin{theorem}[Leibniz公式]\label{theorem:Leibniz公式}
$(f(x)g(x))^{(n)} = \sum_{k = 0}^{n} C_{n}^{k}f^{(n - k)}(x)g^{(k)}(x).$
\end{theorem}

\begin{example}
设\(f(x)\)定义在\([0,1]\)中且\(\lim_{x\rightarrow0^{+}}f\left(x\left(\frac{1}{x}-\left[\frac{1}{x}\right]\right)\right)=0\)，证明：\(\lim_{x\rightarrow0^{+}}f(x)=0\)。
\end{example}
\begin{note}
将极限定义中的$\varepsilon,\delta$适当地替换成$\frac{1}{n},\frac{1}{N}$往往更方便我们分析问题和书写过程.
\end{note}
\begin{proof}

用\(\{x\}\)表示\(x\)的小数部分，则\(x\left(\frac{1}{x}-\left[\frac{1}{x}\right]\right)=x\left\{\frac{1}{x}\right\}\)。

对任意\(\varepsilon>0\)，依据极限定义，存在\(\delta>0\)使得任意\(x\in(0,\delta)\)都有\(\left|f\left(x\left\{\frac{1}{x}\right\}\right)\right|<\varepsilon\)。

取充分大的正整数\(N\)使得\(\frac{1}{N}<\delta\)，则任意\(x\in\left(\frac{1}{N + 1},\frac{1}{N}\right)\)都有\(\left|f\left(x\left\{\frac{1}{x}\right\}\right)\right|<\varepsilon\)。

考虑函数\(x\left\{\frac{1}{x}\right\}\)在区间\(\left(\frac{1}{N + 1},\frac{1}{N}\right)\)中的值域，也就是连续函数
\[g(u)=\frac{u - [u]}{u}=\frac{u - N}{u},u\in(N,N + 1)\]
的值域，考虑端点处的极限可知\(g(u)\)的值域是\(\left(0,\frac{1}{N + 1}\right)\),且严格单调递增.所以对任意\(y\in\left(0,\frac{1}{N + 1}\right)\)，都存在\(x\in\left(\frac{1}{N + 1},\frac{1}{N}\right)\subset(0,\delta)\)使得$\frac{1}{x}=g^{-1}(y)\in(N,N+1)$,即\(y =g(\frac{1}{x})= x\left\{\frac{1}{x}\right\}\)，故\(|f(y)|=\left|f\left(x\left\{\frac{1}{x}\right\}\right)\right|<\varepsilon\).

也就是说，任意\(\varepsilon>0\)，存在正整数\(N\)，使得任意\(y\in\left(0,\frac{1}{N + 1}\right)\)，都有\(|f(y)|<\varepsilon\)，结论得证。

\end{proof}

\begin{example}
设 $f$ 在 $x_0$ 可微且
$$
\alpha_n < x_0 < \beta_n,\ n\in\mathbb{N},\ \lim_{n\to\infty}\alpha_n = \lim_{n\to\infty}\beta_n = x_0,
$$
证明:
$$
\lim_{n\to\infty}\frac{f(\beta_n)-f(\alpha_n)}{\beta_n-\alpha_n}=f'(x_0).
$$
\end{example}
\begin{proof}

不失一般性, 假设 $f(x_0)=x_0=0$, 否则用 $f(x+x_0)-f(x_0)$ 代替 $f(x)$, $\beta_n-x_0$ 代替 $\beta_n$, $\alpha_n-x_0$ 代替 $\alpha_n$ 即可.

现在
\begin{align*}
\lim_{n\to\infty}\frac{f(\beta_n)-f(\alpha_n)}{\beta_n-\alpha_n}
&= \lim_{n\to\infty}\left(\frac{\beta_n}{\beta_n-\alpha_n}\frac{f(\beta_n)}{\beta_n}-\frac{\alpha_n}{\beta_n-\alpha_n}\frac{f(\alpha_n)}{\alpha_n}\right) \\
&= \lim_{n\to\infty}\left(\frac{1}{1-\frac{\alpha_n}{\beta_n}}\frac{f(\beta_n)}{\beta_n}-\frac{\frac{\alpha_n}{\beta_n}}{1-\frac{\alpha_n}{\beta_n}}\frac{f(\alpha_n)}{\alpha_n}\right).
\end{align*}

我们设 $\lim_{k\to\infty}\frac{\alpha_{n_k}}{\beta_{n_k}}=c\in\mathbb{R}$, 则
\begin{align*}
\lim_{k\to\infty}\left(\frac{1}{1-\frac{\alpha_{n_k}}{\beta_{n_k}}}\frac{f(\beta_{n_k})}{\beta_{n_k}}-\frac{\frac{\alpha_{n_k}}{\beta_{n_k}}}{1-\frac{\alpha_{n_k}}{\beta_{n_k}}}\frac{f(\alpha_{n_k})}{\alpha_{n_k}}\right)
= \frac{1}{1-c}f'(x_0)-\frac{c}{1-c}f'(x_0)=f'(x_0).
\end{align*}

设 $\lim_{k\to\infty}\frac{\alpha_{n_k}}{\beta_{n_k}}=\infty$, 则
\begin{align*}
\lim_{k\to\infty}\left(\frac{1}{1-\frac{\alpha_{n_k}}{\beta_{n_k}}}\frac{f(\beta_{n_k})}{\beta_{n_k}}-\frac{\frac{\alpha_{n_k}}{\beta_{n_k}}}{1-\frac{\alpha_{n_k}}{\beta_{n_k}}}\frac{f(\alpha_{n_k})}{\alpha_{n_k}}\right)
= 0\cdot f'(x_0)-(-1)\cdot f'(x_0)=f'(x_0).
\end{align*}
再由\hyperref[proposition:子列极限命题]{子列极限命题}可知,结论成立.

\end{proof}

\begin{example}
设$f(x)\in C(\mathbb{R})$，且有
\begin{align*}
\lim\limits_{h\to 0^+}\frac{f(x+2h)-f(x+h)}{h}=0,\quad \forall x\in \mathbb{R}.
\end{align*}
\begin{enumerate}[(1)]
\item $f(x)$的右导数处处存在.
\item $f(x)$为常值函数.
\end{enumerate}
\end{example}
\begin{proof}
\begin{enumerate}[(1)]
\item 任取$x\in \mathbb{R}$。由题设知对$\forall \varepsilon>0,\exists \delta>0$，使得对$\forall h\in (0,\delta)$，有
\begin{align*}
\left| f\left(x+2h\right)-f\left(x+h\right) \right|<\varepsilon h.
\end{align*}
因为对$\forall k\in \mathbb{N}$，有$\frac{h}{2^{k+1}}\in \left(0,\delta\right)$，所以
\begin{align*}
\left| f\left(x+\frac{h}{2^k}\right)-f\left(x+\frac{h}{2^{k+1}}\right) \right|<\frac{\varepsilon h}{2^k},\quad \forall k\in \mathbb{N}.
\end{align*}
由$f\in C\left(\mathbb{R}\right)$知$\lim\limits_{n\to \infty}\left| f\left(x+\frac{h}{2^n}\right)-f\left(x\right) \right|=0$。从而对$\forall h\in \left(0,\delta\right),n>N$，就有
\begin{align*}
\left| f\left(x+h\right)-f\left(x\right) \right|\leqslant \sum\limits_{k=1}^n\left| f\left(x+\frac{h}{2^{k-1}}\right)-f\left(x+\frac{h}{2^k}\right) \right|+\left| f\left(x+\frac{h}{2^n}\right)-f\left(x\right) \right|\leqslant \sum\limits_{k=1}^n\frac{\varepsilon h}{2^k}+\left| f\left(x+\frac{h}{2^n}\right)-f\left(x\right) \right|.
\end{align*}
令$n\to \infty$得对$\forall h\in \left(0,\delta\right)$，有
\begin{align*}
\left| f\left(x+h\right)-f\left(x\right) \right|\leqslant \varepsilon h\iff \left| \frac{f\left(x+h\right)-f\left(x\right)}{h} \right|\leqslant \varepsilon.
\end{align*}
故$f_{+}'\left(x\right)=\lim\limits_{h\to 0^+}\frac{f\left(x+h\right)-f\left(x\right)}{h}=0$。再由$x$的任意性知$f$的右导数处处为$0$。

\item 任取$a,b\in \mathbb{R}$且$a<b$。由(1)知$f_{+}'\left(x\right)=0\left( \forall x\in \mathbb{R} \right)$，故对$\forall \varepsilon>0,\exists \delta>0$，使得
\begin{align}\label{eq:num31abc}
\left| f\left(x+h\right)-f\left(x\right) \right|\leqslant \varepsilon h.
\end{align}
对闭区间$\left[a,b\right]$做$n\triangleq \lfloor \frac{1}{\delta} \rfloor +1$等分，记分点为
\[a=x_1<x_2<\cdots <x_n<x_{n+1}=b.\]
则
\[x_{i+1}-x_i<\delta ,\quad i=1,2,\cdots ,n.\]
于是由\eqref{eq:num31abc}式知
\begin{align*}
\left| f\left(b\right)-f\left(a\right) \right|\leqslant \sum\limits_{k=1}^n\left| f\left(x_{i+1}\right)-f\left(x_i\right) \right|<\sum\limits_{k=1}^n\varepsilon \left(x_{i+1}-x_i\right)=\varepsilon \left(b-a\right).
\end{align*}
令$\varepsilon \to 0^+$得$f\left(a\right)=f\left(b\right)$。再由$a,b$的任意性知$f\left(x\right)$在$\mathbb{R}$上恒为常数。
\end{enumerate}
\end{proof}




\end{document}